% -----------------------------------------------------------------------------
% Textbook storytelling scaffolding for The Principles of Diffusion Models
% Suggested placement: immediately after 
% Theorem System
% The following boxes are provided:
%   Definition:     \defn 
%   Theorem:        \thm 
%   Lemma:          \lem
%   Corollary:      \cor
%   Proposition:    \prop   
%   Claim:          \clm
%   Fact:           \fact
%   Proof:          \pf
%   例：        \ex
%   Remark:         \rmk (sentence), \rmkb (block)
% Suffix
%   r:              Allow Theorem/Definition to be referenced, e.g. thmr
%   p:              Add a short proof block for Lemma, Corollary, Proposition or Claim, e.g. lemp
%                   For theorems, use \pf for proof blocks

% ======= Real examples : 

% \defn{Definition Name}{
%     A defintion.
% }

% \thmr{Theorem Name}{mybigthm}{
%     A theorem.
% }

% \lem{Lemma Name}{
%     A lemma.
% }

% \fact{
%     A fact.
% }

% \cor{
%     A corollary.
% }

% \prop{
%     A proposition.
% }

% \clmp{}{
%     A claim.
% }{
%     A reference to Theorem~\ref{thm:mybigthm}
% }

% \pf{
%    proof.

% \ex{
%     some examples. 
% }

% \rmk{
%     Some remark.
% }

% \rmkb{
%     Some more remark.
% }
\usepackage{xcolor}

% Define custom colors
\definecolor{defscol}{HTML}{ecd8d7} %For definitions
\definecolor{asumscol}{HTML}{ecd8d7} %For Assumptions

\definecolor{rmkscol}{HTML}{313160} %For remarks
\definecolor{exmscol}{HTML}{e04b52} %For examples

\definecolor{lemscol}{HTML}{2c3943} %For Lemmes
\definecolor{thmscol}{RGB}{230,120,20}
\definecolor{prpscol}{HTML}{8A9A5B}
\definecolor{corscol}{HTML}{dfd9fd} %For corrolaries

\definecolor{clmscol}{HTML}{165c58} %For claims
\definecolor{facscol}{HTML}{28a8a1} %For facts


% ============================
% Definition
% ============================
\newtcbtheorem[number within=section]{mydefinition}{定义}
{
    enhanced,
    frame hidden,
    titlerule=0mm,
    toptitle=1mm,
    bottomtitle=1mm,
    fonttitle=\sffamily\bfseries,
    coltitle=black,
    colbacktitle=defscol!50!white,
    colback=defscol!20!white,
}{defn}

\NewDocumentCommand{\defn}{m+m}{
    \begin{mydefinition}{#1}{}
        #2
    \end{mydefinition}
}

\NewDocumentCommand{\defnr}{mm+m}{
    \begin{mydefinition}{#1}{#2}
        #3
    \end{mydefinition}
}

% ============================
% Problem Setup
% ============================
\newtcbtheorem[number within=section]{myproblem}{问题设定}
{
    enhanced,
    frame hidden,
    titlerule=0mm,
    toptitle=1mm,
    bottomtitle=1mm,
    fonttitle=\sffamily\bfseries,
    coltitle=black,
    colbacktitle=defscol!50!white,
    colback=defscol!20!white,
}{defn}

\NewDocumentCommand{\probs}{m+m}{
    \begin{myproblem}{#1}{}
        #2
    \end{myproblem}
}

\NewDocumentCommand{\problems}{mm+m}{
    \begin{myproblem}{#1}{#2}
        #3
    \end{myproblem}
}

% ============================
% Message
% ============================

% Define a counter that resets at each \section
\newcounter{msgcounter}[section]
\renewcommand{\themsgcounter}{\thesection.\arabic{msgcounter}}

% Define your custom message box (no need to change this part)
\newtcolorbox{custommessagebox}[2][]{
    enhanced,
    coltitle=black,
    colbacktitle=defscol!50!white,
    colback=defscol!20!white,
    fonttitle=\sffamily\bfseries,
    title=#1,        % title passed in
    toptitle=1mm,
    bottomtitle=1mm,
    titlerule=0mm,
    frame hidden,
    #2
}

% Define the \msg command with auto-numbered title based on section
\NewDocumentCommand{\msg}{mm+m}{
    \refstepcounter{msgcounter} % increments the counter and allows referencing
    \begin{custommessagebox}[#1~\themsgcounter:~\textbfs{#2}]{}
        #3
    \end{custommessagebox}
}


% ============================
% Research Problem
% ============================
\newtcbtheorem{myresearchproblem}{研究问题}
{
    enhanced,
    frame hidden,
    titlerule=0mm,
    toptitle=1mm,
    bottomtitle=1mm,
    fonttitle=\sffamily\bfseries,
    coltitle=black,
    colbacktitle=defscol!40!white,
    colback=defscol!20!white,
}{defn}

\NewDocumentCommand{\probr}{m+m}{
    \begin{myresearchproblem}{#1}{}
        #2
    \end{myresearchproblem}
}

\NewDocumentCommand{\problemr}{mm+m}{
    \begin{myresearchproblem}{#1}{#2}
        #3
    \end{myresearchproblem}
}


% ============================
% Assumption
% ============================
\newtcbtheorem[use counter from=mydefinition]{myassumption}{假设}
{
    enhanced,
    frame hidden,
    titlerule=0mm,
    toptitle=1mm,
    bottomtitle=1mm,
    fonttitle=\sffamily\bfseries,
    coltitle=black,
    colbacktitle=asumscol!40!white,
    colback=asumscol!20!white,
}{asum}

\NewDocumentCommand{\asum}{m+m}{
    \begin{myassumption}{#1}{}
        #2
    \end{myassumption}
}

\NewDocumentCommand{\asumr}{mm+m}{
    \begin{myassumption}{#1}{#2}
        #3
    \end{myassumption}
}

% ============================
% Theorem
% ============================

% Custom theorem counter that resets every chapter (not section!)
% Number theorems within sections
\numberwithin{equation}{section}

% Set up a custom counter for theorems
\newcounter{mytheorem}[section]
\renewcommand{\themytheorem}{\Roman{chapter}.\thesection.\arabic{mytheorem}}

% Define the theorem environment with custom label
\newtcbtheorem[use counter=mytheorem, number within=section]{mytheorem}{定理}{
    enhanced,
    frame hidden,
    titlerule=0.0mm,
    fonttitle=\sffamily\bfseries,
    coltitle=black,
    colbacktitle=thmscol!20!white,
    colback=thmscol!8!white,
}{thm}

% Proof environment (page-break friendly)
\makeatletter
\newenvironment{thmpf}{%
  \par\noindent{\it \textbfs{定理的证明。}}\par
  % Key bits: breakable + enhanced jigsaw + allowdisplaybreaks
  \tcolorbox[
    blanker,
    breakable,
    enhanced jigsaw,
    parbox=false,
    colback=white,
    left=5mm,
    before upper={\parindent15pt\relax\allowdisplaybreaks},
    after skip=10pt,
    borderline west={1mm}{0pt}{thmscol!20!white}
  ]%
}{%
  \textcolor{thmscol!18!white}{\hbox{}\nobreak\hfill$\blacksquare$}%
  \endtcolorbox
}
\makeatother


% Combined theorem macro
\newcommand{\thm}[3]{%
    \begin{mytheorem}{#1}{#2}%
        #3%
    \end{mytheorem}%
     \par
}


% Combined theorem+proof macro
\newcommand{\thmp}[4]{%
    \begin{mytheorem}{#1}{#2}%
        #3%
    \end{mytheorem}%
     \par
    \begin{thmpf}%
        #4%
    \end{thmpf}%
    
}


% ============================
% Lemma (breakable-ready)
% ============================
\newtcbtheorem[
  use counter=mytheorem,
  number within=section
]{mylemma}{引理}{
  enhanced,
  % Add 'breakable' here ONLY if you want the *lemma box* to split across pages:
  % breakable,
  frame hidden,
  titlerule=0mm,
  toptitle=1mm,
  bottomtitle=1mm,
  fonttitle=\sffamily\bfseries,
  coltitle=black,
  colbacktitle=corscol!50!white,
  colback=corscol!20!white,
}{lem}

% \lem{title}{label}{content}
%             ^^^^^  ^^^^^^^
% Make 'content' long: +m
\NewDocumentCommand{\lem}{m m +m}{%
  \begin{mylemma}[label=#2]{#1}{}%
    #3%
  \end{mylemma}%
}

% Proof environment for lemma (kept breakable)
\newenvironment{lempf}{%
  \par\noindent{\it \textbfs{引理的证明。}}\par
  \tcolorbox[
    blanker,breakable,left=5mm,parbox=false,
    before upper={\parindent15pt},
    after skip=10pt,
    colback=white,
    borderline west={1mm}{0pt}{corscol!50!white}
  ]%
}{%
  \textcolor{corscol!50!white}{\hbox{}\nobreak\hfill$\blacksquare$}%
  \endtcolorbox
}

% \lemp{title}{label}{content}{proof}
%                             ^^^^^  make both body args long: +m +m
\NewDocumentCommand{\lemp}{m m +m +m}{%
  \begin{mylemma}[label=#2]{#1}{}%
    #3%
  \end{mylemma}%
  \par
  \begin{lempf}%
    #4%
  \end{lempf}%
}

% ============================
% Corollary
% ============================

\newtcbtheorem[use counter from=mydefinition]{mycorollary}{推论}
{
    enhanced,
    frame hidden,
    titlerule=0mm,
    toptitle=1mm,
    bottomtitle=1mm,
    fonttitle=\sffamily\bfseries,
    coltitle=black,
    colbacktitle=corscol!40!white,
    colback=corscol!20!white,
}{cor}

\NewDocumentCommand{\cor}{+m}{
    \begin{mycorollary}{}{}
        #1
    \end{mycorollary}
}

% Proof environment — page-break friendly
\makeatletter
\newenvironment{corpf}{%
  \par\noindent{\it \textbfs{推论的证明。}}\par
  \tcolorbox[
    blanker,
    breakable,          % allow the box to split across pages
    enhanced jigsaw,    % robust page breaking (esp. with math)
    colback=white,
    parbox=false,
    left=5mm,
    before upper={\parindent15pt\allowdisplaybreaks}, % let long displays break
    after skip=10pt,
    borderline west={1mm}{0pt}{corscol!40!white}
  ]%
}{%
  \textcolor{corscol!40!white}{\hbox{}\nobreak\hfill$\blacksquare$}%
  \endtcolorbox
}
\makeatother

\NewDocumentCommand{\corp}{m+m+m}{
    \begin{mycorollary}{{{#1}}}{}
        #2
    \end{mycorollary}

    \begin{corpf}
     {   #3}
    \end{corpf}
}


% Named proposition without proof
\NewDocumentCommand{\cornp}{m+m}{ 
    \begin{mycorollary}{{{#1}}}{}
        #2
    \end{mycorollary}
}


% ============================ % 
% Proposition % 
% ============================
% Define proposition environment with custom label and style
\newtcbtheorem[use counter=mytheorem, number within=section]{myproposition}{命题}
{ 
  enhanced,
  frame hidden,
  titlerule=0mm,
  toptitle=1mm,
  bottomtitle=1mm,
  fonttitle=\sffamily\bfseries,
  coltitle=black,
  colbacktitle=prpscol!18!white,
  colback=prpscol!8!white,
}
{prop}

% Proposition with label and no proof
\NewDocumentCommand{\proppnp}{m m m}{%
  \begin{myproposition}[label=#2]{#1}{}%
    #3%
  \end{myproposition}%
}

% Proof environment for propositions
\newenvironment{proppf}
{ 
  {\noindent{\it \textbfs{命题的证明。}}} 
  \tcolorbox[blanker,breakable,left=5mm,parbox=false, colback=white,
    before upper={\parindent15pt}, after skip=10pt, 
    borderline west={1mm}{0pt}{prpscol!18!white}] 
}
{ 
  \textcolor{prpscol!18!white}{\hbox{}\nobreak\hfill$\blacksquare$} 
  \endtcolorbox 
}

% Proposition with label and proof
\NewDocumentCommand{\proppp}{m m m m}{%
  \begin{myproposition}[label=#2]{#1}{}%
    #3%
  \end{myproposition}%
  \begin{proppf}%
    #4%
  \end{proppf}%
}

% Add this after loading the cleveref package


% ============================
% Claim
% ============================

\newtcbtheorem[use counter from=mydefinition]{myclaim}{断言}
{
    enhanced,
    frame hidden,
    titlerule=0mm,
    toptitle=1mm,
    bottomtitle=1mm,
    fonttitle=\sffamily\bfseries,
    coltitle=black,
    colbacktitle=clmscol!40!white,
    colback=clmscol!20!white,
}{clm}

\NewDocumentCommand{\clm}{m+m}{
    \begin{myclaim*}{#1}{}
        #2
    \end{myclaim*}
}

\newenvironment{clmpf}{
	{\noindent{\it \textbfs{断言的证明。}}}
	\tcolorbox[blanker,breakable,left=5mm,parbox=false,
    before upper={\parindent15pt},
    after skip=10pt,
	borderline west={1mm}{0pt}{clmscol!40!white}]
}{
    \textcolor{clmscol!40!white}{\hbox{}\nobreak\hfill$\blacksquare$} 
    \endtcolorbox
}

\NewDocumentCommand{\clmp}{m+m+m}{
    \begin{myclaim*}{#1}{}
        #2
    \end{myclaim*}

    \begin{clmpf}
        #3
    \end{clmpf}
}

% ============================
% Fact
% ============================

\newtcbtheorem[use counter from=mydefinition]{myfact}{事实}
{
    enhanced,
    frame hidden,
    titlerule=0mm,
    toptitle=1mm,
    bottomtitle=1mm,
    fonttitle=\sffamily\bfseries,
    coltitle=black,
    colbacktitle=facscol!40!white,
    colback=facscol!20!white,
}{fact}

\NewDocumentCommand{\fact}{+m}{
    \begin{myfact}{}{}
        #1
    \end{myfact}
}

% ============================
% Proof
% ============================

\NewDocumentCommand{\pf}{+m}{
    \begin{proof}
        [\noindent\textbfs{证明。}]
        #1
    \end{proof}
}

% ============================
% Example
% ============================


\newenvironment{myexample}{
    \tcolorbox[blanker,breakable,left=5mm,parbox=false,
    colback=white,
    before upper={\parindent15pt},
    after skip=10pt,
	borderline west={1mm}{0pt}{clmscol!25!white}]
}{
    \textcolor{clmscol!25!white}{\hbox{}\nobreak\hfill$\blacksquare$} 
    \endtcolorbox
}

\NewDocumentCommand{\exm}{m+m}{\vspace{0.5cm}
{\noindent{\textbfs{例： #1 }}}
    \begin{myexample}
	
        #2
    \end{myexample}
}


% ============================
% Remark
% ============================


\newenvironment{myremark}{
    \tcolorbox[blanker,breakable,left=5mm,parbox=false,
    colback=white,
    before upper={\parindent15pt},
    after skip=10pt,
	borderline west={1mm}{0pt}{rmkscol!20!white}]
}{
    \textcolor{rmkscol!20!white}{\hbox{}} 
    \endtcolorbox
}

\NewDocumentCommand{\rmkb}{+m}{\vspace{0.5cm}{\par\noindent{\textbfs{注。}}}
    \begin{myremark}
        #1
    \end{myremark}
}


\tcbuselibrary{theorems,skins,hooks}
\newtcbtheorem[number within=section]{Theorem}{定理}
{%
	enhanced,
	breakable,
	colback = mytheorembg,
	frame hidden,
	boxrule = 0sp,
	borderline west = {2pt}{0pt}{mytheoremfr},
	sharp corners,
	detach title,
	before upper = \tcbtitle\par\smallskip,
	coltitle = mytheoremfr,
	fonttitle = \sffamily\bfseries\sffamily,
	description font = \mdseries,
	separator sign none,
	segmentation style={solid, mytheoremfr},
}
{th}


% % Old styles hh 

% %--------------------------------------------------------------------------
% % 		THEOREMES STYLE
% %--------------------------------------------------------------------------

% %-------		DEFINITION 		-------	
% \newcounter{defo}[chapter]
% \newenvironment{defi}[1]{\refstepcounter{defo} 
% \begin{tcolorbox}[colback=yellow!20!white,colframe=yellow!15!black,title= \textbfs{Définition \thechapter \ $\blacklozenge$ \thedefo \ | #1}]}{\end{tcolorbox}}

% %-------		THEOREME 		-------	
% \newcounter{th}[chapter]
% \newenvironment{thm}[1]{\refstepcounter{th}
% \begin{tcolorbox}[colback=mycolor!10,colframe=mycolor!10!black!80,title=\textbfs{ Théorème \thechapter \ $\blacklozenge$ \theth \ | #1}]}{\end{tcolorbox}}

% %-------		PROPOSITION 		-------	
% \newcounter{prop}[chapter]
% \newenvironment{propt}[1]{\refstepcounter{prop}
% \begin{tcolorbox}[colback=mycolor!5,colframe=mycolor!10!linkscolor!80 ,title=\textbfs{ Proposition \thechapter \ $\blacklozenge$ \theprop \ | #1}]}{\end{tcolorbox}}

% %-------		COROLLAIRE		-------
% \newcounter{cor}[chapter]
% \newenvironment{corr}[1]{\refstepcounter{cor}
% \begin{tcolorbox}[colback=mycolor!2,colframe=mycolor!10!linkscolor!40 ,title= Corolaire \thechapter \ $\blacklozenge$ \thecor \ | #1]}{\end{tcolorbox}}

% %-------		LEMME			-------
% \newcounter{lem}[chapter]
% \newenvironment{lemme}[1]{\refstepcounter{lem}
% \begin{tcolorbox}[colback=mycolor!10!blue!2,colframe=mycolor!50!blue!30,title=\textbfs{Lemme \thechapter \ $\blacklozenge$ \thelem \ | #1}]}{\end{tcolorbox}}

% %-------		METHODE 			-------
% \newcounter{met}[chapter]
% \newenvironment{meth}[1]{\refstepcounter{met}
% \begin{tcolorbox}
% [enhanced jigsaw,breakable,pad at break*=1mm,
%  colback=red!20!white,boxrule=0pt,frame hidden,
%  borderline west={1.5mm}{-2mm}{red}] \color{red}
% {\textbfs{Méthode \thechapter \ $\blacklozenge$ \themet \ | #1} } \color{black} \\ } {\end{tcolorbox}}

% %-------		REMARQUE 			-------
% \newcommand{\NB}[1]{
% \ \\
% \begin{tabular}{p{0.05\textwidth}p{0.80\textwidth}}
% \hline
% \vspace{-0.1cm} \includegraphics[scale=0.03]{./system/IDEA.png} & 	#1\\
% \hline
% \end{tabular}
% \ 
% \newline \ \newline
%  }

% %-------		EXEMPLE  		-------
% \newenvironment{exm}{ \begin{tcolorbox}
% [enhanced jigsaw,breakable,pad at break*=1mm,
%  colback=cyan!20!white,boxrule=0pt,frame hidden,
%  borderline west={1.5mm}{-2mm}{cyan}] \color{cyan}
% {\textbfs{Exemple} } \color{black} \\ } {\end{tcolorbox}}


 % ============================
% Unnumbered Message
% ============================

\NewDocumentCommand{\msgu}{m m +m}{
    \begin{custommessagebox}[
        \IfBlankTF{#2}
            {#1}
            {#1:~\textbfs{#2}}
    ]{}
        #3
    \end{custommessagebox}
} in main_v2.tex.
% The definitions use packages already loaded by the manuscript: xcolor,
% tcolorbox, and xparse/NewDocumentCommand.
% -----------------------------------------------------------------------------

% A deliberately small semantic palette. These colors should retain the same
% meaning throughout the book and should also guide new figures.
\definecolor{StoryNavy}{HTML}{23415C}       % chapter/section navigation
\definecolor{StoryBlue}{HTML}{3E789B}       % learned/local model object
\definecolor{StoryBluePale}{HTML}{EEF5F8}
\definecolor{StoryGreen}{HTML}{4F7C63}      % tractable target / takeaway
\definecolor{StoryGreenPale}{HTML}{EEF6F0}
\definecolor{StoryPurple}{HTML}{6D5B8C}     % finite-time map / approximation
\definecolor{StoryPurplePale}{HTML}{F3F0F7}
\definecolor{StoryWarm}{HTML}{9A6A3A}       % caveat / design tension
\definecolor{StoryWarmPale}{HTML}{FAF4EC}
\definecolor{StoryGray}{HTML}{5C6670}       % data/reference/background object
\definecolor{StoryGrayPale}{HTML}{F5F6F7}

% ----------------------------------------------------------------------------
% Chapter guide: one box at the beginning of every chapter.
% Arguments:
%   #1 central question
%   #2 mental model
%   #3 learned object
%   #4 chapter path (ordinary prose or an enumerate/itemize environment)
% ----------------------------------------------------------------------------


\NewDocumentCommand{\chapterguide}{m+m+m+m}{%
  \begin{storychapterguidebox}
  \noindent\textcolor{StoryNavy}{\textbfs{核心问题。}} #1\par\smallskip
  \noindent\textcolor{StoryNavy}{\textbfs{直观图景。}} #2\par\smallskip
  \noindent\textcolor{StoryNavy}{\textbfs{学习什么？}} #3\par\smallskip
  \noindent\textcolor{StoryNavy}{\textbfs{本章路线。}} #4
  \end{storychapterguidebox}
}

% ----------------------------------------------------------------------------
% Lightweight section navigation. Use this at the beginning of every major
% section; it is intentionally not a box, to avoid visual clutter.
% ----------------------------------------------------------------------------
\NewDocumentCommand{\whythissection}{m}{%
  \par\smallskip
  \noindent\textcolor{StoryNavy}{\textbfs{本节的作用。}} #1
  \par\smallskip
}

% ----------------------------------------------------------------------------
% Equation interpretation. Place immediately after a load-bearing equation or
% derivation. It should explain the statement in words, not repeat the symbols.
% ----------------------------------------------------------------------------
\NewDocumentCommand{\equationmeaning}{m}{%
  \par\smallskip
  \noindent\textcolor{StoryBlue}{\textbfs{公式的含义。}} #1
  \par\smallskip
}

% ----------------------------------------------------------------------------
% 要点。 Use at the end of a major section. Keep to one paragraph.
% ----------------------------------------------------------------------------
\newtcolorbox{storytakeawaybox}{
  enhanced,
  breakable,
  frame hidden,
  colback=StoryGreenPale,
  borderline west={1.2mm}{0pt}{StoryGreen},
  left=3.5mm,right=3.5mm,top=2mm,bottom=2mm,
  before skip=8pt,after skip=10pt
}

\NewDocumentCommand{\takeaway}{m}{%
  \begin{storytakeawaybox}
  \noindent\textcolor{StoryGreen}{\textbfs{要点。}} #1
  \end{storytakeawaybox}
}

% ----------------------------------------------------------------------------
% 证明思路。 The optional argument is the appendix location. The main text
% should contain the conceptual mechanism; the complete algebra belongs in the
% proof appendix unless it is shorter than roughly five lines.
% Usage: \proofidea[\Cref{app:proof-foo}]{...}
% ----------------------------------------------------------------------------
\NewDocumentCommand{\proofidea}{o +m}{%
  \par\smallskip
  \noindent\textcolor{StoryPurple}{\textbfs{证明思路。}} #2%
  \IfNoValueF{#1}{\space\textit{完整细节： #1.}}%
  \par\smallskip
}

% ----------------------------------------------------------------------------
% Running example. Reuse the same example and visual grammar across chapters.
% Suggested default: a two-dimensional, three-mode data distribution and a
% Gaussian reference distribution.
% ----------------------------------------------------------------------------
\newtcolorbox{runningexample}[1]{
  enhanced,
  breakable,
  frame hidden,
  colback=StoryGrayPale,
  borderline west={1.2mm}{0pt}{StoryGray},
  title={\textbfs{贯穿示例： #1}},
  fonttitle=\sffamily\bfseries,
  coltitle=StoryNavy,
  colbacktitle=StoryGrayPale,
  titlerule=0mm,
  left=3.5mm,right=3.5mm,top=2mm,bottom=2mm,
  before skip=8pt,after skip=10pt
}

% ----------------------------------------------------------------------------
% Method profile. Use the same five questions for every solver, guidance method,
% distillation method, and flow-map method. This symmetry prevents paper-by-
% paper exposition from reading like a catalog.
% Arguments:
%   #1 method name
%   #2 represented/learned object
%   #3 training or approximation signal
%   #4 generation/sampling mechanism
%   #5 main strength and main limitation
% ----------------------------------------------------------------------------
\newtcolorbox{storymethodbox}[1]{
  enhanced,
  breakable,
  frame hidden,
  colback=StoryWarmPale,
  borderline west={1.2mm}{0pt}{StoryWarm},
  title={\textbfs{方法解析： #1}},
  fonttitle=\sffamily\bfseries,
  coltitle=StoryWarm,
  colbacktitle=StoryWarmPale,
  titlerule=0mm,
  left=3.5mm,right=3.5mm,top=2mm,bottom=2mm,
  before skip=8pt,after skip=10pt
}

\NewDocumentCommand{\methodlens}{m+m+m+m+m}{%
  \begin{storymethodbox}{#1}
  \noindent\textbfs{对象。} #2\par\smallskip
  \noindent\textbfs{学习信号。} #3\par\smallskip
  \noindent\textbfs{生成机制。} #4\par\smallskip
  \noindent\textbfs{优势与局限。} #5
  \end{storymethodbox}
}

% ----------------------------------------------------------------------------
% Chapter close. Use the same five-field syntax in Chapters 1--11. Chapter 12
% may use a book-level variant. Each argument should usually be one sentence.
% ----------------------------------------------------------------------------
\newtcolorbox{storychapterclosebox}{
  enhanced,
  breakable,
  frame hidden,
  colback=StoryPurplePale,
  borderline west={1.4mm}{0pt}{StoryPurple},
  left=4mm,right=4mm,top=2.5mm,bottom=2.5mm,
  before skip=10pt,after skip=8pt
}

\NewDocumentCommand{\chapterclose}{m+m+m+m+m}{%
  \begin{storychapterclosebox}
  \noindent\textcolor{StoryPurple}{\textbfs{已回答的问题。}} #1\par\smallskip
  \noindent\textcolor{StoryPurple}{\textbfs{学习的对象。}} #2\par\smallskip
  \noindent\textcolor{StoryPurple}{\textbfs{可计算的信号。}} #3\par\smallskip
  \noindent\textcolor{StoryPurple}{\textbfs{生成过程。}} #4\par\smallskip
  \noindent\textcolor{StoryPurple}{\textbfs{下一步。}} #5
  \end{storychapterclosebox}
}

% ----------------------------------------------------------------------------
% Optional reading marker. Use at the beginning of an optional section to state
% exactly what a first-time reader may skip and what conceptual fact remains.
% ----------------------------------------------------------------------------
\NewDocumentCommand{\optionalroute}{m+m}{%
  \par\smallskip
  \noindent\textcolor{StoryWarm}{\textbfs{选读路径。}} #1
  \space\textbfs{请记住：} #2
  \par\smallskip
}

% ----------------------------------------------------------------------------
% Recommended semantic color contract for newly revised figures:
%   gray   : fixed data/reference distribution or background geometry
%   blue   : learned local object (score, denoiser, velocity, classifier signal)
%   green  : tractable conditional target or observed supervision
%   red    : inaccessible marginal/oracle target
%   purple : numerical approximation or learned finite-time flow map
% Do not reuse a color for a different mathematical role in another chapter.
% -----------------------------------------------------------------------------
